%% ====================================================================
\section{Conclusiones y recomendaciones}
\label{sec:conclusiones}
%% ====================================================================

Este TFM formula la coordinación distribuida de coaliciones multi-AMR como un problema
de juego poblacional con cierre físico de capacidad. La tesis central es que una flota
industrial no debe decidir solo qué tarea es atractiva, sino qué coalición es capaz de
servir una carga bajo restricciones de quorum, comunicación, energía, contacto,
seguridad y esfuerzo actuador. Smith-QR aporta una solución explícita para ese problema:
mantiene la interpretabilidad de las dinámicas de Smith y añade mecanismos de memoria,
compromiso, patrullaje y cierre local de quorum cuando la conectividad cae por debajo de
lo que exige el equilibrio distribuido ideal.

Las conclusiones principales son las siguientes.

\begin{enumerate}
  \item El caso estático homogéneo queda explicado por la regla de
  \emph{water-filling}. El protocolo mínimo reproduce pendiente $1.0000$ y $R^2=1$,
  lo que confirma que Smith reparte masa hacia tareas deficitarias hasta igualar payoffs
  marginales.

  \item El \emph{clearing} entero es necesario porque una asignación fluida puede ser
  matemáticamente coherente y físicamente inútil. Al proyectar preferencias continuas a
  coaliciones enteras se reducen desplazamientos desperdiciados y se evitan sub-quorums
  incapaces de transportar una carga.

  \item El precio temporal no debe interpretarse como una mejora universal. En equilibrio
  estático puede perturbar un reparto ya resuelto por Smith; su valor aparece cuando hay
  llegadas, \emph{deadlines} y déficit persistente. Así, el precio se entiende como señal
  dinámica de urgencia, no como corrección permanente del equilibrio.

  \item Smith original depende de conectividad suficiente. En el régimen R3 los robots
  pueden perseguir tareas compatibles con su información local, pero quedar repartidos en
  grupos que no cierran contacto físico. Smith-QR corrige esa debilidad al convertir la
  preferencia poblacional en cierre local de quorum.

  \item En la campaña R3 de 20 semillas, Smith-QR mejora a Smith al pasar de $0.1038$ a
  $0.2893$ de captura media de recompensa. El delta pareado es $+0.1855$, con IC95
  $[0.1434,0.2276]$, y el número medio de entregas pasa de $3.25$ a $8.25$. También
  supera en captura a greedy local y al proxy MARL de referencia en esta campaña, aunque la distancia
  desperdiciada no mejora de forma concluyente. La conclusión correcta es multiobjetivo:
  Smith-QR rescata la dinámica Smith y mejora captura bajo R3, pero debe evaluarse junto
  a coste de movimiento, energía y batería.

  \item La exploración logit/mutation no resuelve por sí sola la parálisis de R3. El
  problema no es únicamente falta de ruido exploratorio, sino falta de cierre físico y
  comunicación suficiente. Por eso la solución robusta debe modificar la regla de quorum,
  no solo añadir aleatoriedad a la selección de tarea.

  \item MARL aporta una referencia algorítmica relevante cuando se entrena y se evalúa
  como baseline real. La política MARL-CTDE lineal aprende pesos compartidos con
  entrenamiento centralizado y ejecución descentralizada; mejora al proxy en el régimen
  de escasez, pero no domina a Smith-QR globalmente. El actor neuronal CTDE añade una
  capa no lineal residual: mejora el objetivo de entrenamiento, supera al proxy en
  escasez y reduce distancia desperdiciada frente al actor lineal en ese régimen, aunque
  tampoco domina a Smith-QR. Esa conclusión es útil: Smith-QR no se compara sólo contra
  una heurística débil, sino contra competidores aprendidos cuya ventaja y límite quedan
  medidos con $H_0/H_1$.

  \item La extensión de capacidad efectiva transforma el criterio
  $|\mathcal C_k|\ge n_k$ en una condición física
  $S_k(\mathcal C_k,t)\ge D_k(t)$. Esta formulación permite incorporar masa, torque,
  fricción, geometría, obstáculos, batería, formación y saturación de actuadores. La
  coalición deja de ser un conteo de robots y pasa a ser un cierre de restricciones
  mecánicas.

  \item La escalabilidad se separa en dos niveles. El nivel defendible en esta memoria es
  la aproximación determinista de la dinámica de Smith: bajo comunicación global y
  escalado $\beta_N=\tilde\beta/N$, la flota finita se concentra alrededor de la ODE con
  error $O_p(N^{-1/2})$. El nivel MFG queda como extensión algorítmica futura para campos
  escalares, densidad espacial y peajes dinámicos.

  \item El transporte cooperativo amorfo completa la teoría para cargas rígidas. La flota
  se modela como densidad en espacio de fase; la carga, como sistema port-Hamiltoniano; y
  el contacto, como integral de wrench sobre el contorno. Con ello, asignación, formación,
  seguridad y empuje se escriben como partes de una misma energía acoplada.

  \item El framework integrado conecta economía, batería, formación, inercia, torques,
  campo medio y métricas en el funcional $\mathcal F_{\mathrm{int}}$. Su valor está en
  ordenar el problema por capas verificables: potencial de Smith para la decisión,
  pasividad para la carga, energía geométrica para la formación, retornabilidad para la
  batería, aproximación determinista para el puente finito--ODE y MFG como ruta futura de
  escalabilidad espacial. La nueva lectura por cadena de residuos
  \eqref{eq:cadena_motor_integrado}--\eqref{eq:residuos_motor_integrado} impide que esas
  piezas queden como ideas separadas: cada capa consume una señal de déficit y entrega
  una magnitud auditable a la capa siguiente.

  \item La robustez del framework exige abandonar tres idealizaciones del modelo base:
  Laplaciano casi estático, capacidad escalar e información simétrica perfecta. La
  guardia anti-\emph{chattering} del DAC, la capacidad en espacio wrench y la regla
  marginal inspirada en VCG convierten esas debilidades en condiciones de diseño
  medibles.

  \item Los acoplamientos biofísicos ofrecen una lectura adicional del reclutamiento y
  la formación. La inferencia activa recluta robots cuando existe error físico de
  predicción en la velocidad de la carga; Kuramoto permite sincronizar fases de contacto
  para que el \emph{caging} aparezca como patrón emergente alrededor del objeto. Ambos
  mecanismos complementan Smith-QR sin destruir su trazabilidad.

  \item El control vectorial ejecutable convierte la teoría en un bucle implementable
  para CoppeliaSim. Los contornos circulares y superelípticos cubren cargas circulares,
  cuadradas y rectangulares; la firma fuerza--torque del contorno conecta MFG con wrench
  físico; el residual de wrench reemplaza el conteo escalar por déficit físico;
  Smith--Kuramoto ajusta fases de contacto; la rigidez local mantiene el macro-sólido; el
  QP CBF filtra seguridad; y el mapeo a uniciclo entrega comandos $(v_i,\omega_i)$
  verificables por ciclo.

  \item Los gates G3 y G4 cierran la parte algebraicamente auditable del motor avanzado.
  G3 demuestra que el quorum escalar puede aceptar una coalición físicamente inútil para
  producir torque. G4 valida que geometría, mercado, peaje predictivo, batería y
  port-Hamiltoniano usan señales compatibles: el círculo radial no genera torque, el
  cuadrado y el rectángulo sí producen firmas de torque, el peaje mueve la fase de
  contacto, la batería ordena payoffs y el balance energético respeta pasividad.
\end{enumerate}

Como recomendación técnica, la validación debe mantenerse multiobjetivo. No basta medir
entregas o distancia; también deben compararse energía total, potencia de rueda,
saturación de torque, distancia mínima, violaciones de seguridad, vida útil de batería,
tiempo de control y robustez frente a comunicación degradada. Un método que aumenta
throughput a costa de saturar actuadores o consumir batería de retorno no resuelve el
problema industrial completo.

La línea de trabajo más prometedora es una arquitectura híbrida: Smith-QR como núcleo
interpretable de asignación y quorum; MARL-CTDE lineal y actor neuronal CTDE como módulos
complementarios para exploración, ajuste de parámetros o selección de hiperpolíticas; capacidad wrench para decidir qué robots son físicamente útiles;
inferencia activa para reclutar por error físico de carga; sincronización de fases para
formación espontánea; control vectorial con CBF y uniciclo para ejecutar la coalición
finita; y campo medio para controlar congestión cuando la flota crece. Esa
combinación mantiene ecuaciones defendibles y, al mismo tiempo, deja espacio para
políticas aprendidas donde el modelado analítico sea demasiado costoso.
